% ============================================================
%  cap05/5.6_contraste_eficiencia.tex
% ============================================================
\subsection{Contraste de Eficiencia y Seguridad en Arquitecturas Multi-Tenant}

El n\'ucleo arquitect\'onico de \textit{Democra} orbita alrededor de un dise\~no \textit{Multi-Tenant} de tipo \textit{Pool} (Esquema Compartido, Datos Compartidos), asegurado a nivel de tupla mediante pol\'iticas de \textit{Row-Level Security} (RLS) inyectadas directamente en el kernel del motor relacional PostgreSQL 16. La disyuntiva acad\'emica cl\'asica al implementar este tipo de mallas compartidas radica en la tensi\'on entre el m\'aximo aprovechamiento de hardware (ahorro de costos) y la degradaci\'on latente del rendimiento computacional (Sobrecarga de Filtrado).

\subsubsection{Sobrecarga Computacional (RLS Overhead)}
Nuestros resultados de laboratorio demostraron que la habilitaci\'on del RLS nativo, apoyado sobre JWT y par\'ametros de configuraci\'on local de sesi\'on (\texttt{request.jwt.claims}), impuso una degradaci\'on cronom\'etrica de apenas un $4.22\%$ (una fluctuaci\'on de milisegundos pr\'acticamente imperceptible para el ruteo de red HTTP convencional). Por el contrario, la emulaci\'on de un filtrado cl\'asico por software (Capa de Aplicaci\'on) dispar\'o la sobrecarga computacional a un inasumible $501.40\%$.

Esta abismal brecha de eficiencia se alinea sin\'ergicamente y corrobora de forma concluyente los hallazgos postulados por \textbf{Red Hat Research (2023)} \citep{redhat2023}. Dicho estudio europeo advirti\'o expl\'icitamente que la resoluci\'on de permisos iterativos fuera del kernel de la base de datos anula la capacidad del Planificador de Consultas (\textit{Query Planner}) para emplear \'indices eficientemente. En sinton\'ia con \textit{Red Hat}, \textit{Democra} demostr\'o que al cristalizar la identificaci\'on del \textit{Tenant} directamente como una macro de base de datos indexada en \'arboles B (\textit{B-Tree}), el motor omite la inspecci\'on de p\'aginas de memoria ajenas. En consecuencia, la decisi\'on de aislar los dominios operativos de m\'ultiples ONGs utilizando RLS nativo no solo fue exitosa en mitigar fugas de datos (alcanzando un ratio de intercepci\'on defensiva del $100\%$), sino que result\'o en una topolog\'ia econ\'omicamente escalable que elude el aprovisionamiento excesivo (\textit{Overprovisioning}).
